% ============================================================
% Chapter 3: Problem Formulation
% Thesis: Solving the FSRCPSP Using Hybrid ML Approaches
% Author: Souvik Chakraborty
% ============================================================
\chapter{Problem Formulation}
\label{chap:formulation}

This chapter formally defines the optimisation problem addressed
in this thesis. Section~\ref{sec:problem_desc} describes the
Flexible Stochastic Resource-Constrained Project Scheduling
Problem (FSRCPSP) and its defining characteristics.
Section~\ref{sec:notation} introduces the sets, parameters, and
decision variables. Section~\ref{sec:math_model} presents the
complete SAA-FSRCPSP formulation, with each constraint explained
individually. Finally, Section~\ref{sec:model_summary} summarises
the model structure and discusses its computational properties.

% ------------------------------------------------------------
\section{Problem Description}
\label{sec:problem_desc}
% ------------------------------------------------------------

We consider a project consisting of $n$ non-preemptive activities
that must be scheduled subject to precedence and resource
constraints. The project is represented as an Activity-on-Node
network $\mathcal{N} = (V, E)$, where
$V = \{0, 1, \ldots, n, n+1\}$ denotes the set of activities and
$E \subseteq V \times V$ the set of directed precedence arcs
\parencite{Pritsker1969}. Node set $V$ consists of $n$ real
activities $\{1, \ldots, n\}$ and two fictitious activities: a
source node $0$ denoting the unique start and a sink node $n+1$
denoting the unique end of the project. An arc $(i, j) \in E$
enforces a finish-to-start precedence relationship, indicating
that activity~$i$ must be completed before activity~$j$ can be
started.

The set of renewable resources is denoted
$K = \{1, \ldots, |K|\}$, where each resource $k \in K$ has a
fixed capacity $R_k$ available in every discrete time period
$t \in T = \{0, 1, \ldots, T_{\max}\}$. The planning horizon
$T_{\max}$ is determined by a longest-path calculation on the
network. For each activity $i \in V$, the earliest and latest
feasible start times $ES_i$ and $LS_i$ are computed via forward
and backward passes over the network respectively
\parencite{Kelley1963}. Set
$\overline{W}_i = \{ES_i, ES_i + 1, \ldots, LS_i\}$ defines
the possible start times of activity~$i$.

The FSRCPSP extends the classical RCPSP in two directions
\parencite{Mendl2024}. First, each activity~$i$ has a
\emph{flexible resource profile}: rather than consuming a fixed
amount per time period, the resource allocation $r_{ikt} \geq 0$
may vary within predefined lower and upper resource limits
$\underline{r}_{ik} \geq 0$ and
$\overline{r}_{ik} \leq \max_{t \in T}\{R_k\}$ across the
activity's duration, provided that the total workload is exactly
covered \parencite{Naber2014}. Second, the workload $w_{ik}$ of
activity~$i$ on resource~$k$ is \emph{stochastic}. A
deterministic workload $w^d_{ik}$ defines the base workload for
each resource and the ratio between the different resource
requirements per activity. A stochastic scaling factor $f_i > 0$
for activity~$i$ is drawn from a predefined probability
distribution $P$. The stochastic workload is then calculated as
$w_{ik} = w^d_{ik} \cdot f_i$, preserving the resource demand
ratios across resources \parencite{Mendl2024}. All activities are
assumed to be non-preemptive, meaning each resource $k \in K$
must be allocated for the entire duration $d_i$ of activity~$i$
and cannot be interrupted.

To incorporate uncertainty into the model, we introduce random
variables for the stochastic workloads and approximate them
through a discrete scenario set $S = \{1, \ldots, |S|\}$, where
each scenario $\pi \in S$ corresponds to one specific realisation
of the workload distributions. The resulting stochastic model is
transformed into an equivalent deterministic model using a Sample
Average Approximation (SAA) approach with chance-constraints, as
in \textcite{Bruni2011}, \textcite{Lamas2016}, and \textcite{MendlRieck2024}. The
construction of the chance-constraints depends on the effects
that different workloads have on the decision variables. In our
problem, the workloads particularly affect the durations $d_i$ of
activities and the resource allocations $r_{ikt}$. Thus, the
conditions for maintaining precedence constraints and for
maintaining resource capacities are formulated as
chance-constrained conditions \parencite{Sallam2021}.

The objective is to determine a proactive, robust baseline
schedule $S^B$ that minimises the project duration. Following
\textcite{Lamas2016} and \textcite{MendlRieck2024}, we measure robustness through a
\emph{confidence level} $(1 - \alpha)$, set by the project
manager, which specifies the probability that the schedule remains
feasible under any realisation of the stochastic workloads.
Constraints~(\ref{eq:precedence}) and~(\ref{eq:resource_cap})
in the original chance-constrained formulation require
$\Pr(\cdot) \geq 1 - \alpha$ over the true distribution
\parencite{Lamas2016}. Since this probabilistic model cannot
be solved directly, the SAA replaces the distributional guarantee
by an empirical one over the scenario set~$S$: a schedule is
accepted if it is feasible in at least a fraction
$(1 - \varepsilon)$ of the sampled scenarios, where
$\varepsilon$ denotes the sample-based violation tolerance
(constraint~\ref{eq:saa_cc}). Setting
$\varepsilon \leq \alpha$ ensures that the SAA solution satisfies
the original chance constraint with high confidence
\parencite{Lamas2016}. In this thesis, we use
$\varepsilon = 0.1$ with $|S| = 10$ scenarios, so at most
$\lfloor 10 \cdot 0.1 \rfloor = 1$ scenario may be violated.

% ------------------------------------------------------------
\section{Sets, Parameters and Decision Variables}
\label{sec:notation}
% ------------------------------------------------------------

Table~\ref{tab:sets_params} summarises the sets and parameters
of the model.

\begin{table}[htbp]
\centering
\caption{Sets and parameters of the SAA-FSRCPSP}
\label{tab:sets_params}
\renewcommand{\arraystretch}{1.1}
\begin{tabular}{lp{9cm}}
\hline
\textbf{Symbol} & \textbf{Description} \\
\hline
\multicolumn{2}{l}{\textit{Sets}} \\[2pt]
$V$       & All activities including source $0$ and sink $n+1$ \\[4pt]
$E$       & Set of directed precedence arcs \\[4pt]
$K$       & Set of renewable resources $\{1,\ldots,|K|\}$ \\[4pt]
$T$       & Discrete time periods $\{0,1,\ldots,T_{\max}\}$ \\[4pt]
$S$       & Set of SAA scenarios $\{1,\ldots,|S|\}$ \\[6pt]
\hline
\multicolumn{2}{l}{\textit{Parameters}} \\[2pt]
$ES_i$    & Earliest feasible start time of activity $i$ \\[4pt]
$LS_i$    & Latest feasible start time of activity $i$ \\[4pt]
$\overline{W}_i$
          & $[ES_i, LS_i]$: feasible time window of activity $i$ \\[4pt]
$R_k$     & Capacity of resource $k \in K$ per time period \\[4pt]
$\underline{r}_{ik}$
          & Minimum resource allocation of $i$ on $k$ per period \\[4pt]
$\overline{r}_{ik}$
          & Maximum resource allocation of $i$ on $k$ per period \\[4pt]
$w_{ik}^{\pi}$
          & Workload of activity $i$ on resource $k$ in scenario $\pi$ \\[4pt]
$\varepsilon$
          & Violation tolerance ($\varepsilon = 0.1$ in this thesis) \\
\hline
\end{tabular}
\end{table}

To model the SAA-FSRCPSP, we use two groups of binary
decision variables. Following \textcite{Pritsker1969}, we specify
$x^B_{it} \in \{0, 1\}$ and its continuous relaxation
$x^C_{it} \in [0, 1]$, $i \in V$, $t \in T$, which together form
the effective start variable
$x^{\mathrm{eff}}_{it} = x^B_{it} + x^C_{it}$. The variable
$x^{\mathrm{eff}}_{it}$ indicates whether activity~$i$ is started
at time~$t$. The splitting into binary and continuous components
is central to the Relax-and-Fix decomposition described in
Chapter~\ref{chap:baseline}: in each subproblem, a subset of
activities is treated as binary (using $x^B_{it}$) while the
remaining activities are relaxed to continuous (using $x^C_{it}$)
\parencite{pochet2006, Helber2010}.

The binary variable $y_{it}^{\pi} \in \{0, 1\}$, $i \in V$,
$t \in T$, $\pi \in S$, indicates that activity~$i$ has not yet
completed at time~$t$ in scenario~$\pi$. The real-valued variables
$r_{ikt}^{\pi} \geq 0$ and $d_i^{\pi} \geq 0$, $i \in V$,
$k \in K$, $t \in T$, $\pi \in S$, represent the resource
allocation of activity~$i$ on resource~$k$ at time~$t$ in
scenario~$\pi$ and the duration of activity~$i$ in scenario~$\pi$,
respectively. A scenario-level binary $\rho_{\pi} \in \{0, 1\}$
takes the value 1 if any precedence or resource constraint is
violated in scenario~$\pi$. The decision variables are summarised in
Table~\ref{tab:dec_vars}.

\begin{table}[htbp]
\centering
\caption{Decision variables of the SAA-FSRCPSP}
\label{tab:dec_vars}
\renewcommand{\arraystretch}{1.4}
\begin{tabular}{lp{9cm}}
\hline
\textbf{Variable} & \textbf{Description} \\
\hline
$x^B_{it} \in \{0,1\}$
    & 1 if activity $i$ starts at time $t$ (binary component) \\[4pt]
$x^C_{it} \in [0,1]$
    & Continuous relaxation of $x^B_{it}$, used in R\&F
      subproblems \\[4pt]
$x^{\mathrm{eff}}_{it} \in [0,1]$
    & Effective start variable:
      $x^{\mathrm{eff}}_{it} = x^B_{it} + x^C_{it}$ \\[4pt]
$r_{ikt}^{\pi} \geq 0$
    & Resource allocation of $i$ on $k$ at time $t$ in
      scenario $\pi$ \\[4pt]
$d_i^{\pi} \geq 0$
    & Realised duration of activity $i$ in scenario $\pi$ \\[4pt]
$y_{it}^{\pi} \in \{0,1\}$
    & 1 if activity $i$ has not yet completed at time $t$ in
      scenario $\pi$ \\[4pt]
$\rho_{\pi} \in \{0,1\}$
    & 1 if scenario $\pi$ is violated, 0 otherwise \\
\hline
\end{tabular}
\end{table}

% ------------------------------------------------------------
\section{Mathematical Formulation}
\label{sec:math_model}
% ------------------------------------------------------------

The SAA-FSRCPSP can be formulated as follows
\parencite{Mendl2024}:

\subsection*{Objective Function}

\begin{equation}
    \min \quad \sum_{t \in T} t \cdot x^{\mathrm{eff}}_{n+1,\, t}
    \tag{1}
    \label{eq:obj}
\end{equation}

The objective~(\ref{eq:obj}) is to minimise the project duration.
Since the sink activity $n+1$ has zero duration, its start time
equals the project makespan $C_{\max} = S_{n+1}$.

\subsection*{Constraint (2): Each Activity Starts Exactly Once}

\begin{equation}
    \sum_{t \in \overline{W}_i} x^{\mathrm{eff}}_{it} = 1
    \qquad \forall\, i \in V
    \tag{2}
    \label{eq:start_once}
\end{equation}

Constraints~(\ref{eq:start_once}) ensure that each activity~$i$
is started exactly once within its feasible time window
$\overline{W}_i$. This enforces the non-preemptive assumption: an
activity cannot be split or restarted once it has begun.

\subsection*{Constraint (3): Precedence Constraints}

\begin{equation}
    \sum_{t \in \overline{W}_j} t \cdot x^{\mathrm{eff}}_{jt}
    - \sum_{t \in \overline{W}_i} t \cdot x^{\mathrm{eff}}_{it}
    \geq d_i^{\pi}
    - \left(|T| + \max_{k \in K}
      \left\lceil \frac{w_{ik}^{\pi}}{\underline{r}_{ik}}
      \right\rceil\right) \cdot \rho_{\pi}
    \quad \forall\, (i,j) \in E,\;
    i \in V \setminus \{0, n{+}1\},\;
    \pi \in S
    \tag{3}
    \label{eq:precedence}
\end{equation}

\subsection*{Constraint (4): Resource Capacity}

\begin{equation}
    \sum_{i \in V} r_{ikt}^{\pi}
    - \left(\sum_{i \in V} \overline{r}_{ik} - R_k\right)
      \cdot \rho_{\pi}
    \leq R_k
    \quad \forall\, k \in K,\;
    t \in T,\;
    \pi \in S
    \tag{4}
    \label{eq:resource_cap}
\end{equation}

Constraints~(\ref{eq:precedence}) and~(\ref{eq:resource_cap})
check if a solution is feasible under the scenario $\pi \in S$
with regard to the precedence constraints and resource
capacities. The big-$M$ coefficients used in these conditions are
already suitably small due to the maximum functions. In
constraint~(\ref{eq:precedence}), the term
$|T| + \max_{k} \lceil w_{ik}^{\pi} / \underline{r}_{ik} \rceil$
equals the horizon length plus the maximum possible duration of
activity~$i$ when it is executed at its minimum rate
$\underline{r}_{ik}$. In constraint~(\ref{eq:resource_cap}), the
coefficient $\sum_{i} \overline{r}_{ik} - R_k$ represents the
maximum possible resource excess when all activities
simultaneously consume their upper rate limits. The binary
decision variables $\rho_{\pi} \in \{0, 1\}$ take the value 1 if
any violation occurs in scenario $\pi \in S$, meaning that either
the precedence constraints and/or the resource constraints are not
satisfied.

\subsection*{Constraint (5): Workload Coverage}

\begin{equation}
    \sum_{t=ES_i}^{LS_i} r_{ikt}^{\pi}
    = w_{ik}^{\pi}
    \quad \forall\, \pi \in S,\;
    k \in K,\;
    i \in V
    \tag{5}
    \label{eq:workload_cover}
\end{equation}

Equations~(\ref{eq:workload_cover}) specify that the total amount
of allocated resources $r_{ikt}^{\pi}$ must exactly cover the
required workload $w_{ik}^{\pi}$ of activity~$i$ for resource~$k$
in scenario~$\pi$. This is the defining constraint of the
flexible resource profile: the workload is fixed by the scenario,
but how it is distributed over time is a decision variable
\parencite{Naber2014}.

\subsection*{Constraint (6): Big-M Linking of $y_{it}^{\pi}$}

\begin{equation}
    y_{it}^{\pi}
    \leq \sum_{k \in K} \sum_{\tau \geq t} r_{ik\tau}^{\pi}
    \leq \left(\sum_{k \in K} w_{ik}^{\pi}\right)
         \cdot y_{it}^{\pi}
    \quad \forall\, i \in V,\;
    t \in T,\;
    \pi \in S
    \tag{6}
    \label{eq:bigm}
\end{equation}

Constraints~(\ref{eq:bigm}) link the allocation variables to the
completion indicator: resources can only be allocated to
activity~$i$ at time~$t$ if it has not yet been completed. Once the activity is finished
($y_{it}^{\pi} = 0$), no further resources can be assigned to it.
The left inequality forces $y_{it}^{\pi} = 0$ when no resource
allocation remains, and the right inequality forces
$y_{it}^{\pi} = 1$ when any allocation remains.

\subsection*{Constraint (7): Duration Identification}

\begin{equation}
    \sum_{t \in T}
    \left(
        y_{it}^{\pi}
        + \sum_{\tau \leq t} x^{\mathrm{eff}}_{i\tau} - 1
    \right)
    = d_i^{\pi}
    \quad \forall\, i \in V \setminus \{0, n{+}1\},\;
    \pi \in S
    \tag{7}
    \label{eq:duration}
\end{equation}

By connecting the decision variables $x^{\mathrm{eff}}_{it}$ and
$y_{it}^{\pi}$ in condition~(\ref{eq:duration}), it can be
determined whether an activity~$i$ is in progress at time~$t$ or
not. The number of periods in which resources are assigned to
activity~$i$ corresponds to its realised duration $d_i^{\pi}$.

\subsection*{Constraint (8): Flexible Resource Allocation Bounds}

\begin{equation}
    \underline{r}_{ik}
    \left(
        y_{it}^{\pi}
        + \sum_{\tau \leq t} x^{\mathrm{eff}}_{i\tau} - 1
    \right)
    \leq r_{ikt}^{\pi}
    \leq \overline{r}_{ik}
    \left(
        y_{it}^{\pi}
        + \sum_{\tau \leq t} x^{\mathrm{eff}}_{i\tau} - 1
    \right)
    \quad \forall\, i \in V,\;
    k \in K,\;
    t \in T,\;
    \pi \in S
    \tag{8}
    \label{eq:alloc_bounds}
\end{equation}

Constraints~(\ref{eq:alloc_bounds}) ensure that the allocated
resources $r_{ikt}^{\pi}$ are always within the lower
$\underline{r}_{ik}$ and upper limits $\overline{r}_{ik}$ of
activity~$i$ for resource~$k$. When activity~$i$ is active at
time~$t$ in scenario~$\pi$, the activity-progress term
$(y_{it}^{\pi} + \sum_{\tau \leq t} x^{\mathrm{eff}}_{i\tau} - 1)$
equals 1 and the allocation must lie within
$[\underline{r}_{ik},\, \overline{r}_{ik}]$. When $i$ is not
active, the term equals zero and the constraint forces
$r_{ikt}^{\pi} = 0$.

\subsection*{Constraint (9): SAA Chance Constraint}

\begin{equation}
    \sum_{\pi \in S} \rho_{\pi}
    \leq \lfloor |S| \cdot \varepsilon \rfloor
    \tag{9}
    \label{eq:saa_cc}
\end{equation}

Inequality~(\ref{eq:saa_cc}) limits the total number of violated
scenarios to at most $\lfloor |S| \cdot \varepsilon \rfloor$,
enforcing the confidence level $(1 - \varepsilon)$. With
$|S| = 10$ and $\varepsilon = 0.1$, at most
$\lfloor 10 \cdot 0.1 \rfloor = 1$ scenario may be violated.

\subsection*{Constraints (10)--(12): Auxiliary Consistency
Constraints}

\begin{equation}
    \sum_{t \in T} y_{it}^{\pi}
    \geq d_i^{\pi}
    \quad \forall\, i \in V,\;
    \pi \in S
    \tag{10}
    \label{eq:consistency1}
\end{equation}

\begin{equation}
    y_{it}^{\pi} = 1
    \quad \forall\, i \in V \setminus \{0, n{+}1\},\;
    t \leq ES_i,\;
    \pi \in S
    \tag{11}
    \label{eq:consistency2}
\end{equation}

\begin{equation}
    y_{it}^{\pi} \geq x^{\mathrm{eff}}_{it}
    \quad \forall\, i \in V \setminus \{0, n{+}1\},\;
    t \leq ES_i,\;
    \pi \in S
    \tag{12}
    \label{eq:consistency3}
\end{equation}

These three auxiliary constraints tighten the LP relaxation and
ensure internal consistency of the model.
Constraint~(\ref{eq:consistency1}) requires that the total number
of active periods is at least the realised duration.
Constraint~(\ref{eq:consistency2}) fixes $y_{it}^{\pi} = 1$ for
all periods before $ES_i$: since the activity has not yet started,
it is by definition not yet completed.
Constraint~(\ref{eq:consistency3}) ensures that $y_{it}^{\pi}$ is
at least as large as $x^{\mathrm{eff}}_{it}$ in the same early
periods, maintaining consistency between the start and progress
variables.

% ------------------------------------------------------------
\section{Complete Model Summary}
\label{sec:model_summary}
% ------------------------------------------------------------

The complete SAA-FSRCPSP minimises~(\ref{eq:obj}) subject to
constraints~(\ref{eq:start_once})--(\ref{eq:consistency3}),
with variable domains:
\begin{align*}
    x^B_{it}      &\in \{0,1\}
    && \forall\, i \in V,\; t \in T \\
    x^C_{it}      &\in [0,1]
    && \forall\, i \in V,\; t \in T \\
    r_{ikt}^{\pi} &\geq 0
    && \forall\, i \in V,\; k \in K,\;
       t \in T,\; \pi \in S \\
    d_i^{\pi}     &\geq 0
    && \forall\, i \in V,\; \pi \in S \\
    y_{it}^{\pi}  &\in \{0,1\}
    && \forall\, i \in V,\; t \in T,\;
       \pi \in S \\
    \rho_{\pi}    &\in \{0,1\}
    && \forall\, \pi \in S
\end{align*}

The model is a mixed-integer programme whose total variable
count is
\[
  \underbrace{|V| \cdot |T|}_{\displaystyle x^B}
  + \underbrace{|V| \cdot |T|}_{\displaystyle x^C}
  + \underbrace{|V| \cdot |T|}_{\displaystyle x^{\mathrm{Eff}}}
  + \underbrace{|V| \cdot |T| \cdot |S|}_{\displaystyle y}
  + \underbrace{|S|}_{\displaystyle \rho}
  + \underbrace{|V| \cdot |S|}_{\displaystyle d}
  + \underbrace{|V| \cdot |K| \cdot |T| \cdot |S|}_{\displaystyle r}
\]
with the dominant term being the $r$-variables, indexed over
$V \times K \times T \times S$.
Table~\ref{tab:model_size} reports the resulting model
dimensions for the benchmark instances used in this thesis.

\begin{table}[htbp]
\centering
\caption{Model dimensions for benchmark instances
         ($|K| = 4$, $|S| = 10$ scenarios), computed analytically
         from the variable and constraint formulas above.
         Ranges reflect the variation in the planning
         horizon~$|T|$ across the instances per class.}
\label{tab:model_size}
\renewcommand{\arraystretch}{1.3}
\begin{tabular}{lcrr}
\hline
\textbf{Instance class} &
$|T|$ &
\textbf{Variables} &
\textbf{Constraints} \\
\hline
$n = 10$, $|K| = 4$
    & 70--92
    & 44{,}650--58{,}642
    & 88{,}233--115{,}777 \\
$n = 30$, $|K| = 4$
    & 130--199
    & 220{,}810--337{,}834
    & 426{,}973--652{,}741 \\
$n = 60$, $|K| = 4$
    & 269--409
    & 884{,}564--1{,}344{,}604
    & 1{,}698{,}381--2{,}580{,}661 \\
\hline
\end{tabular}
\end{table}

\noindent The $r$-variables alone account for approximately
75\,\% of all decision variables. As Table~\ref{tab:model_size}
shows, the model grows roughly $6\times$ from $n=10$ to $n=30$
and a further $4\times$ from $n=30$ to $n=60$, making exact
solution of larger instances intractable without decomposition.
This motivates the Relax-and-Fix approach
presented in Chapters~\ref{chap:baseline}
and~\ref{chap:grouprf}, where the model is solved iteratively
through a sequence of smaller subproblems, and the scalability
extensions of Section~\ref{sec:scalability}, which further reduce
the time and scenario dimensions of each subproblem.
